\documentclass[11pt]{article}
\setlength{\parindent}{0cm}
\usepackage{parskip}
\usepackage[margin=.9in]{geometry}
\usepackage{setspace}
\setstretch{1.1}
\usepackage{titling}
\setlength{\droptitle}{-4em} 
\usepackage{titlesec}
\titlespacing*{\section}{0pt}{10pt plus 2pt minus 2pt}{5pt plus 1pt minus 1pt}
\renewcommand{\maketitlehookd}{\vspace{-2em}}
\usepackage{enumitem}
\usepackage{longtable}
\usepackage{array}
\usepackage{hyperref}

\title{DAEN 300: Data Engineering Coding Experience I}
\author{Industrial and Systems Engineering, Texas A\&M University}
\date{Fall 2027}

\begin{document}

\maketitle

\section*{Course Information}

\begin{itemize}[leftmargin=*,nosep]
%    \item \textbf{Course:} DAEN 300 -- Data Engineering Coding Experience I
    \item \textbf{Credit Hours:} 1, \textbf{Contact Hours:} 0 lecture hours, 3 laboratory hours per week
    \item \textbf{Prerequisites:} Grade of C or better in ECEN 360;
    concurrent enrollment in DAEN 321 and DAEN 331;
    junior or senior classification.
    \item \textbf{Format:} Hands-on computational laboratory, \textbf{Programming Language:} Python
    \item \textbf{Instructor:} Alexander Abuabara, \textbf{Office hours:} Wednesdays 1-3 PM, ETB 4013\\
    \url{https://engineering.tamu.edu/industrial/profiles/abuabara-alexander.html}
\end{itemize}

\section*{Course Description}

Application of computational tools to model and solve data
engineering problems primarily involving machine learning
and optimization techniques.

This laboratory course develops practical programming
skills for solving data engineering problems. Students
prepare and analyze data, implement machine learning models,
apply optimization techniques such as loss minimization
and gradient descent, and solve basic constrained optimization
problems. Activities emphasize interpretation
of results and reinforce concepts from concurrent courses.

%This laboratory course emphasizes practical programming
%skills for solving data engineering problems.
%Students use Python and computational libraries to prepare
%and analyze data, implement machine learning algorithms,
%formulate optimization problems, and interpret computational
%results. Activities reinforce concepts introduced in
%concurrent data engineering courses.

%This laboratory course emphasizes practical programming
%skills for solving data engineering problems.
%Students use Python and computational libraries to prepare
%and analyze data, implement and evaluate machine learning
%models, and apply optimization techniques underlying
%machine learning, including loss minimization, gradient
%descent, and model selection. The course also introduces
%constrained optimization using SciPy. Students learn to
%interpret computational results and communicate their
%findings. Laboratory activities reinforce concepts
%introduced in concurrent data engineering courses.

\section*{Course Learning Objectives}

Upon successful completion of this course, students will
be able to:

\begin{enumerate}[label=\arabic*.,nosep]
    \item Develop Python programs to process, analyze,
    and visualize engineering datasets.
    \item Implement and evaluate basic machine learning
    models for prediction and classification.
    \item Formulate and solve constrained optimization
    problems using computational tools.
    \item Interpret computational results, assess model
    performance, and communicate findings.
    \item Develop reproducible and documented computational
    solutions, including collaborative workflows.
\end{enumerate}

%\section*{Alignment with ABET Student Outcomes}
%
%The course supports the following ABET Engineering
%Accreditation Commission student outcomes:
%
%\begin{itemize}[leftmargin=*,nosep]
%    \item \textbf{Outcome 1:} Identify, formulate, and solve
%    complex engineering problems by applying principles of
%    engineering, science, and mathematics.
%    
%    \item \textbf{Outcome 2:} Apply engineering design to
%    produce solutions that meet specified needs with
%    consideration of relevant constraints.
%    
%    \item \textbf{Outcome 3:} Communicate effectively with
%    a range of audiences.
%    
%    \item \textbf{Outcome 5:} Function effectively on a team.
%    
%    \item \textbf{Outcome 6:} Develop and conduct appropriate
%    experimentation, analyze and interpret data, and use
%    engineering judgment to draw conclusions.
%    
%    \item \textbf{Outcome 7:} Acquire and apply new knowledge
%    as needed, using appropriate learning strategies.
%\end{itemize}
%
%\noindent
%The primary outcomes assessed in this course are
%Outcomes 1, 3, and 6.

\section*{Course Topics}

\begin{enumerate}[nosep]
    \item \textbf{Computational Foundations}
    \begin{itemize}[nosep]
        \item Python and Jupyter notebooks
        \item Numerical computing with NumPy
        \item Data manipulation with pandas
        \item Data visualization
        \item Code documentation and version control
    \end{itemize}

    \item \textbf{Machine Learning}
    \begin{itemize}[nosep]
        \item Regression
        \item Classification
        \item Training and testing datasets
        \item Model evaluation and validation
        \item Engineering prediction applications
    \end{itemize}

    \item \textbf{Optimization}
    \begin{itemize}[nosep]
        \item Objective functions and constraints
        \item Linear programming
        \item Numerical optimization
        \item Engineering resource allocation problems
    \end{itemize}

    \item \textbf{Integrated Applications}
    \begin{itemize}[nosep]
        \item Combining prediction and optimization
        \item Reproducible computational workflows
        \item Interpretation and communication of results
    \end{itemize}
\end{enumerate}

\section*{Tentative Course Schedule}

\begin{longtable}{|>{\centering\arraybackslash}p{0.9cm}|p{9cm}|}
\hline
\textbf{Week} & \textbf{Laboratory Topic} \\
\hline
1  & Python environment, Jupyter notebooks, and Git \\ \hline
2  & Numerical computing with NumPy \\ \hline
3  & Data manipulation and cleaning with pandas \\ \hline
4  & Exploratory data analysis and visualization \\ \hline
5  & Linear regression \\ \hline
6  & Classification \\ \hline
7  & Model validation and performance evaluation \\ \hline
8  & Machine learning case study \\ \hline
9  & Optimization formulation and constraints \\ \hline
10 & Linear programming with SciPy \\ \hline
11 & Constrained numerical optimization \\ \hline
12 & Engineering optimization case study \\ \hline
13 & Integrated project development \\ \hline
14 & Integrated project implementation and testing \\ \hline
15 & Project demonstration and presentation \\ \hline
\end{longtable}

\section*{Software and Computational Tools}

The course uses freely available computational tools:

\begin{itemize}[leftmargin=*, nosep]
    \item Python 3, JupyterLab or Google Colab
    \item NumPy and pandas
    \item Matplotlib
    \item scikit-learn
    \item SciPy optimization tools
    \item Git and GitHub
\end{itemize}

\section*{Course Materials}

\textbf{Primary Reference:}

\begin{itemize}[leftmargin=*, nosep]
\item \emph{Data Science: A First Introduction (Python Edition).}
      University of British Columbia.\\
      \url{https://github.com/UBC-DSCI/introduction-to-datascience-python}
\end{itemize}

\textbf{Supplementary Reference:}

\begin{itemize}[leftmargin=*, nosep]
\item \emph{Mathematical Programming and Operations Research.}\\
      \url{https://github.com/open-optimization/open-optimization-or-book}
\item Additional Python notebooks, datasets, and laboratory
      instructions will be provided by the instructor.
\end{itemize}

\section*{Assessment and Grading}

\begin{center}
\begin{tabular}{lr}
\hline
\textbf{Assessment Component} & \textbf{Weight} \\
\hline
Weekly coding laboratories & 60\% \\
Machine learning case study & 15\% \\
Integrated final project & 20\% \\
Participation and code documentation & 5\% \\
\hline
\textbf{Total} & \textbf{100\%} \\
\hline
\end{tabular}
\end{center}

Assignments will be evaluated based on correctness,
appropriate use of computational methods, interpretation
of results, and code clarity and reproducibility.

\section*{Integrated Final Project}

Students will complete a small computational project
combining machine learning and optimization.

The project should include:

\begin{enumerate}[nosep]
    \item Data preparation and exploratory analysis.
    \item Development and evaluation of a predictive model.
    \item Formulation of an optimization problem.
    \item Computational solution and interpretation.
    \item Submission of a documented Jupyter notebook
    and a brief summary of findings.
\end{enumerate}

Possible applications include demand forecasting,
manufacturing planning, and resource allocation.

\section*{Instructional Approach}

The course emphasizes learning through hands-on coding.
Laboratory meetings will consist of short demonstrations,
guided exercises, and independent or collaborative
programming activities.

Students are expected to apply mathematical and engineering
concepts introduced in prerequisite and concurrent courses.

\section*{Course Expectations}

Students are expected to:

\begin{itemize}[leftmargin=*,nosep]
    \item Attend and participate in laboratory sessions.
    \item Complete and submit coding assignments on time.
    \item Write readable and appropriately documented code.
    \item Interpret and explain computational results.
    \item Follow university academic integrity policies.
    \item Acknowledge external code, libraries, and
    AI-assisted contributions as required by the instructor.
\end{itemize}

\end{document}